\chapter{Normas de matrices}\label{cap:normas-matrices}

En el Capítulo~\ref{cap:normas} estudiamos normas en espacios vectoriales
cualesquiera, y en particular las de $\R^{n\times p}$: como ese espacio tiene
dimensión $np$, todo lo de allí se aplica a matrices sin cambiar una coma.
Pero las matrices tienen una operación que los vectores no tienen ---el
producto--- y eso cambia las cosas. Este capítulo trata de las normas de
matrices que \emph{son compatibles con el producto}. Esa compatibilidad es lo
que las vuelve la herramienta central del análisis numérico: es la que
permite acotar el error de un cálculo largo, decidir si un método iterativo
converge y medir cuán confiable es la solución de un sistema lineal.

% ===================================================================
\section{Normas matriciales}
\label{sec:normas-matriciales}

Recordemos la construcción del Capítulo~\ref{cap:normas}. Sea
$\textup{vec}:\R^{n\times p}\rightarrow\R^{np}$ la función que transforma
matrices en vectores apilando sus columnas. Para cada norma $\|\cdot\|$ en
$\R^{np}$ queda definido un funcional $\vertiii{\cdot}$ en $\R^{n\times p}$,
\[
\vertiii{A}=\|\textup{vec}(A)\| ,
\]
y no es difícil mostrar que es una norma (Ejercicio~\ref{ejer-normaiii}). El
caso más usado es la \textbf{norma de Frobenius} de la
Sección~\ref{sec:frobenius},
\[
\|A\|_F=\|\textup{vec}(A)\|_2=\sqrt{\sum_{i=1}^n\sum_{j=1}^p a_{ij}^2}
=\sqrt{\tr(A^TA)} ,
\]
que no es otra cosa que la norma $\|\cdot\|_2$ de $\R^{np}$.

Hasta acá nada distingue a una matriz de un vector largo. La diferencia
aparece cuando se multiplican. Si queremos usar una norma para acotar errores
en una cuenta que encadena productos, lo mínimo que necesitamos es poder
acotar $\|AB\|$ en términos de $\|A\|$ y $\|B\|$. Eso es lo que pedimos.

\begin{defn}[Norma matricial]\label{def:norma-matricial}
Una norma $\|\cdot\|$ definida sobre matrices es una \textbf{norma matricial}
si además de las tres propiedades de norma verifica la
\emph{submultiplicatividad}:
\begin{enumerate}[label=(n\arabic*)]
\setcounter{enumi}{3}
\item $\|AB\|\le\|A\|\,\|B\|$ para todo par de matrices $A$ y $B$ compatible
con el producto.
\end{enumerate}
\end{defn}

Conviene subrayar que la desigualdad es $\le$ y no $=$: pedir la igualdad
sería absurdo, porque $AB$ puede ser la matriz nula sin que lo sean $A$ ni
$B$. Lo que (n4) dice es que \emph{multiplicar no puede agrandar más que el
producto de los tamaños}, y eso alcanza para que, iterando,
$\|A^k\|\le\|A\|^k$.

\begin{prop}\label{prop:frobenius-matricial}
    La norma de Frobenius es una norma matricial.
\end{prop}

\begin{proof}
Sean $A\in\R^{m\times n}$ y $B\in\R^{n\times p}$. La entrada $(i,j)$ de $AB$ es
el producto escalar de la fila $i$ de $A$ con la columna $j$ de $B$,
\[
(AB)_{ij}=\fila_i(A)\,\col_j(B) ,
\]
así que por la desigualdad de Cauchy--Schwarz
\[
(AB)_{ij}^2\ \le\ \|\fila_i(A)\|_2^2\;\|\col_j(B)\|_2^2 .
\]
Ahora sumamos en $i$ y en $j$. El primer factor no depende de $j$ y el segundo
no depende de $i$, de modo que la suma doble se parte en un producto de dos
sumas simples:
\[
\|AB\|_F^2\ \le\ \Bigl(\sum_{i=1}^{m}\|\fila_i(A)\|_2^2\Bigr)
\Bigl(\sum_{j=1}^{p}\|\col_j(B)\|_2^2\Bigr)=\|A\|_F^2\,\|B\|_F^2 ,
\]
porque la norma de Frobenius al cuadrado es tanto la suma de los cuadrados de
las longitudes de las filas como la de las columnas
(Ejercicio~\ref{ej-matrizortogonal}). Tomando raíz cuadrada se completa la
demostración.
\end{proof}

\begin{remark}[Una cota mejor, mirando por columnas]\label{obs:frob-por-columnas}
La fórmula~\eqref{eq:cols} del Capítulo~\ref{cap:algebra-lineal} dice que
$\col_j(AB)=A\,\col_j(B)$: el producto $AB$ no es otra cosa que $A$ aplicada a
cada columna de $B$, una por una. Si medimos cada una de esas columnas con la
norma espectral,
\[
\|AB\|_F^2=\sum_{j=1}^{p}\|A\,\col_j(B)\|_2^2
\ \le\ \|A\|_2^2\sum_{j=1}^{p}\|\col_j(B)\|_2^2=\|A\|_2^2\,\|B\|_F^2 ,
\]
es decir $\|AB\|_F\le\|A\|_2\|B\|_F$. Como $\|A\|_2\le\|A\|_F$, esto es más
fuerte que la proposición anterior. La versión simétrica,
$\|AB\|_F\le\|A\|_F\|B\|_2$, sale mirando por filas
con~\eqref{eq:filas} (Ejercicio~\ref{ej-mixtas}).
\end{remark}

No es cierto que toda norma en un espacio de matrices sea una norma
matricial. Para verlo basta mirar la que se obtiene de la norma del máximo de
$\R^{np}$,
\[
\|A\|_{\max}=\|\textup{vec}(A)\|_\infty=\max_{i,j}|a_{ij}| ,
\]
---el mayor elemento en módulo--- que llamaremos \textbf{norma del máximo}.
Tomando
\[
A=B=\begin{bmatrix}
1 & 1 \\
1 & 1
\end{bmatrix}
\]
resulta $\|AB\|_{\max}=2$ mientras que $\|A\|_{\max}\|B\|_{\max}=1$: la
desigualdad (n4) falla. El Ejercicio~\ref{ej-medica} muestra que
multiplicándola por una constante adecuada se la puede arreglar.

\begin{remark}[Frobenius y los valores singulares]\label{obs:frobenius-sigmas-rec}
La Proposición~\ref{prop:frobenius-sigmas} del Capítulo~\ref{cap:SVD} dice que
\[
\|A\|_F^2=\sigma_1^2+\cdots+\sigma_r^2 ,
\]
y el Corolario~\ref{coro:tres-invariantes} que $\|\cdot\|_F$, $\|\cdot\|_2$ y
$\|\cdot\|_*$ son ortogonalmente invariantes. Que la norma nuclear sea
efectivamente una norma es el Ejercicio~\ref{ej-nuclear}.
\end{remark}

% ===================================================================
\section{Normas inducidas}
\label{sec:normas-inducidas}
Una clase de normas sobre matrices que resultan ser siempre normas matriciales son las que consideran a la matriz $A \in \RR^{n\times p}$ no como un vector reordenado sino como un operador lineal entre $\RR^p$ y $\RR^n$. Dada una norma vectorial $\| \cdot \|$ definida tanto en $\RR^n$ como en $\RR^p$ definimos la norma matricial inducida por $\|\cdot\|$ como
\begin{equation}\label{eq-normainducida}
    \|A\| = \max_{\|\bec x\|\neq 0} \frac{\|A\bec x\|}{\|\bec x\|}.
\end{equation}
Aquí lo estándar es usar la misma notación que para la norma vectorial como para la norma matricial inducida. Antes de demostrar que toda norma inducida es una norma matricial hacemos la siguiente observación que quedará como Ejercicio~\ref{ej-equivalencianormainducida}:
\[\max_{\bec x \neq \bec 0}\frac{\|A\bec x\|}{\|\bec x\|} = \max_{\|\bec x \| = 1}{\|A\bec x\|} = \max_{\|\bec x \| \leq  1}{\|A\bec x\|}.   \]

\begin{thm}\label{T:consistencia}
    Dada una norma vectorial $\| \cdot \|$, la norma definida por~\eqref{eq-normainducida} es una norma matricial.
\end{thm}

\begin{proof}
     Verificaremos que la norma inducida cumple con las tres propiedades de una norma: no negatividad y propiedad de separación, homogeneidad, y desigualdad triangular. Luego, probaremos la submultiplicatividad.



        Por definición, tenemos
        \[
        \|A\| = \max_{\|\bec x\| \neq 0} \frac{\|A\bec x\|}{\|\bec x\|}.
        \]
        Dado que tanto \(\|A\bec x\|\) como \(\|\bec x\|\) son no negativas, el cociente también es no negativo, por lo que \(\|A\| \geq 0\).
Si \(A = 0\), entonces \(A\bec x = 0\) para todo \(\bec x\), lo cual implica que \(\|A\bec x\| = 0\) para todo \(\bec x\), y, por lo tanto, \(\|A\| = 0\).
Por otro lado, si \(\|A\| = 0\), entonces \(\|A\bec x\| = 0\) para todo \(\bec x\) con \(\|\bec x\| \neq 0\), lo cual solo es posible si \(A = 0\). 

 
        Sea ahora \(\alpha \in \mathbb{R}\). Tenemos que
        \[
        \|\alpha A\| = \max_{\|\bec x\| \neq 0} \frac{\|\alpha A \bec x\|}{\|\bec x\|} = \max_{\|\bec x\| \neq 0} \frac{|\alpha| \|A\bec x\|}{\|\bec x\|} = |\alpha| \max_{\|\bec x\| \neq 0} \frac{\|A\bec x\|}{\|\bec x\|} = |\alpha| \|A\|.
        \]
        Esto prueba la propiedad de homogeneidad.


        Para verificar la desigualdad triangular queremos probar que \(\|A + B\| \leq \|A\| + \|B\|\). Por definición de la norma inducida,
        \[
        \|A + B\| = \max_{\|\bec x\| \neq 0} \frac{\|(A + B)\bec x\|}{\|\bec x\|}.
        \]
        Dado que \(\|(A + B)\bec x\| \leq \|A\bec x\| + \|B\bec x\|\) (por la desigualdad triangular de la norma vectorial), se tiene que
        \[
        \frac{\|(A + B)\bec x\|}{\|\bec x\|} \leq \frac{\|A\bec x\|}{\|\bec x\|} + \frac{\|B\bec x\|}{\|\bec x\|}.
        \]
        Tomando el máximo en ambos lados respecto de \(\|\bec x\| \neq 0\), obtenemos
        \[
        \|A + B\| = \max_{\|\bec x\| \neq 0} \frac{\|(A + B)\bec x\|}{\|\bec x\|} \leq \max_{\|  \bec x\| \neq 0} \frac{\|A\bec x\|}{\|\bec x\|} + \max_{\|\bec x\| \neq 0} \frac{\|B\bec x\|}{\|\bec x\|} = \|A\| + \|B\|.
        \]


    Por lo tanto, \(\|A\|\) cumple con las propiedades de no negatividad, homogeneidad y desigualdad triangular, lo cual demuestra que es una norma.
    Finalmente, probaremos que \(\|AB\| \leq \|A\| \|B\|\).

    Por definición,
    \[
    \|AB\| = \max_{\|\bec x\| \neq 0} \frac{\|AB \bec x\|}{\|\bec x\|}.
    \]
    Sea \(\bec y = B\bec x\). Entonces \(\|\bec y\| \leq \|B\| \|\bec x\|\), y
    \[
    \|AB \bec x\| = \|A \bec y\| \leq \|A\| \|\bec y\| \leq \|A\| \|B\| \|\bec x\|.
    \]
    Por lo tanto,
    \[
    \frac{\|AB \bec x\|}{\|\bec x\|} \leq \|A\| \|B\|,
    \]
    y tomando el máximo sobre todos los \(\bec x\) con \(\|\bec x\| \neq 0\), obtenemos que
    \[
    \|AB\| \leq \|A\| \|B\|.
    \]
\end{proof}


La definición~\eqref{eq-normainducida} es un máximo sobre infinitos vectores:
sirve para demostrar, pero no para calcular. Para las tres normas vectoriales
más usadas ---$\|\cdot\|_2$, $\|\cdot\|_1$ y $\|\cdot\|_\infty$--- ese máximo
se puede escribir directamente en términos de las entradas de la matriz.

\begin{thm}\label{teo-l1l2linfty} Sea $A$ una matriz de columnas $\{\bec a_1, \dots, \bec a_n\}$ y filas $\{\bec a_1^\star, \dots \bec a_m^\star\}$. Tenemos que 
    \begin{enumerate}[label=\textup{(\alph*)}]
        \item\label{item-l2} $\|A\|_2 = \sigma_1(A)$.
        \item\label{item-l1} \(\displaystyle\|A\|_1  = \max_{1
        \leq j
        \leq n} \sum_i |a_{i,j}| = \max_{1\leq j\leq n} \|\bec a_j\|_1. \)
        \item\label{item-linfty} \(\displaystyle\|A\|_\infty = \max_{1\leq i \leq m} \sum_j |a_{i,j}| =  \max_{1\leq i\leq m} \|\bec a_i^\star\|_1. \) 
    \end{enumerate}
\end{thm}

\begin{proof}
La parte~\ref{item-l2} ya está hecha: la norma inducida por $\|\cdot\|_2$ es
exactamente la \emph{norma espectral} de la
Definición~\ref{def:norma-espectral}, y el Teorema~\ref{teo:sigma1} dice que
ese máximo vale $\sigma_1(A)$. Dicho de otro modo, la norma espectral que
venimos usando desde el Capítulo~\ref{cap:SVD} es un caso particular de la
construcción de esta sección, y por el Teorema~\ref{T:consistencia} resulta
ser, de paso, una norma matricial.

Ahora, para la norma $\ell_1$, la fórmula~\eqref{eq:mat-vec-col} del
Capítulo~\ref{cap:algebra-lineal} dice que $A\bec x$ es la combinación de las
columnas de $A$ con coeficientes $x_j$. Entonces, por la desigualdad
triangular, si $\bec x = [x_1,\dots,x_n]^T$
\begin{equation*}
    \|A\bec x\|_1 = \left\|\sum_{j=1}^n x_j \bec a_j\right\|_1 \leq \sum_{j=1}^n \| x_j \bec a_j\|_1 = \sum_{j=1}^n |x_j|\| \bec a_j\|_1.
\end{equation*}
Por lo tanto, si $\|x\|_1 = 1$,
\begin{equation*}
     \|A\bec x\|_1 \leq \left(\max_{1\leq j\leq n}\| \bec a_j\|_1\right) \sum_{j=1}^n |x_j|  = \max_{1\leq j\leq n}\| \bec a_j\|_1
\end{equation*}
y
\begin{equation*}
    \|A\|_1 = \max_{\|\bec x\|_1 = 1}\|A\bec x\|_1 \leq \max_{1\leq j\leq n}\| \bec a_j\|_1. 
\end{equation*}

Para mostrar la igualdad, consideremos una matriz $A$ y sea $j_0$ tal que 
\[\max_{1\leq j \leq m}\|\bec a_j\|_1 = \|\bec a_{j_0}\|_1.\]
Si tomamos $\bec x = \bec e_{j_0}$, el vector que vale  uno en la posición $j_0$ y cero en las demás, tenemos que
\[\|A\bec x\|_1 = \|\bec a_{j_0}\| = \max_{1\leq j \leq m}\|\bec a_j\|_1.\]
Como además $\|\bec x\|_1 = 1$ tenemos que 
\begin{equation*}
    \|A\|_1 \geq \max_{1\leq j\leq n}\| \bec a_j\|_1. 
\end{equation*}

La caracterización de la norma $\ell_\infty$ sale del mismo modo, pero mirando
por filas en lugar de por columnas: queda como el
Ejercicio~\ref{ej-linfty}.
\end{proof}

\section{Espectro y radio espectral.}

Dada una matriz $A\in\Rnn$ se define su \emph{espectro} $\sigma(A)$ como el conjunto de autovalores de $A$, así
\[
\sigma(A) = \{ \lambda \in \C : \lambda \text{ es autovalor de }A\}.
\]

El \emph{radio espectral} $\rho(A)$ de $A$ es el máximo módulo de sus autovalores, más precisamente
\[
\rho(A) = \max \{ |\lambda| : \lambda \in \sigma(A) \}.
\]
Por lo dicho anteriormente,
\[
\text{si $A$ es simétric}a \quad\Longrightarrow\quad \rho(A) = \|A\|_2.
\]

El radio espectral juega un papel fundamental en el estudio de métodos iterativos para sistemas lineales, principalmente debido al siguiente teorema:

\begin{thm}[Caracterización del radio espectral]
  \label{T:radio-espectral}
  Sea $A\in\Rnn$, entonces
  \[
  \rho(A) = \inf \{ \|A\| : \|\cdot\| \text{ es una norma inducida} \}.
  \]
\end{thm}


La demostración se apoya en la descomposición de Schur
(Teorema~\ref{teo:schur} del Capítulo~\ref{cap:schur}): toda
$A\in\C^{n\times n}$ se escribe $A=UTU^{*}$ con $U$ unitaria y $T$ triangular
superior cuya diagonal son los autovalores de $A$ contados con multiplicidad.
Eso es cierto aunque $A$ no sea diagonalizable, y si los autovalores de $A$ son
reales se puede tomar $U$ real ortogonal.

\begin{proof}[Demostración del Teorema~\ref{T:radio-espectral}]
  Es fácil ver que $\rho(A) \le \|A\|$ para toda norma inducida $\| \cdot \|$ y cualquier matriz cuadrada $A$ (EJERCICIO). Lo más difícil es ver que dada $A\in\Rnn$, para todo $\varepsilon > 0$ existe una norma inducida $\| \cdot \|$ tal que $\|A\| < \rho(A) + \varepsilon$.

  Sean entonces $A\in\Rnn$ y $\varepsilon > 0$ dados.
  Por el Teorema~\ref{teo:schur} de descomposición de Schur existe $U$ unitaria tal que $UA\bar U^T = T$ con $T$ triangular superior tal que
  \[
  \diag(T) = (\lambda_1, \lambda_2,\dots,\lambda_n),\quad\text{y $\lambda_i$, $i=1,2,\dots,n$ son los autovalores de $A$}.
  \]
  Entonces $T$ es de la forma
  \[
  T = \begin{bmatrix}
    \lambda_1 & t_{12}  & t_{13} & \dots & t_{1n} \\
    0      & \lambda_2 & t_{23} & \dots & t_{2n} \\
    \vdots  & \ddots    & \ddots & \ddots & \vdots  \\
    0       & \dots  & 0      & \lambda_{n-1} & t_{(n-1)n}\\
    0       & \dots  & 0      & 0     & \lambda_n
  \end{bmatrix}.
  \]
  Para $z \in \R$, $z > 0$, definimos
  \[
  D_z = \diag(z,z^2,\dots,z^n) =
  \begin{bmatrix}z & 0   & 0   & \dots & 0 \\
    0 & z^2 & 0   & \dots & 0 \\
    \vdots & \ddots & \ddots & \ddots & \vdots \\
    0 & 0 & \dots & z^{n-1} & 0 \\
    0 & 0 & \dots & 0 & z^n
  \end{bmatrix}
  \]
  Por~\eqref{eq:filas}, multiplicar por $D_z$ a izquierda multiplica la fila
  $i$ por $z^i$; por~\eqref{eq:cols}, multiplicar por $D_z^{-1}$ a derecha
  multiplica la columna $j$ por $z^{-j}$. Haciendo las dos cosas, la entrada
  $(i,j)$ de $T$ queda multiplicada por $z^{i-j}$: la diagonal no se mueve
  ---ahí $i=j$--- y todo lo de arriba de la diagonal, donde $i<j$, se achica.
  Es decir,
  \[
  D_z  U A \bar U^T D_z^{-1} =
  D_z T D_z^{-1} = \begin{bmatrix}
    \lambda_1 & z^{-1} t_{12} & z^{-2} t_{13}  & \dots & z^{-(n-1)} t_{1n} \\
    0      & \lambda_2 & z^{-1}t_{23}  & \dots & z^{-(n-2)}t_{2n} \\
    \vdots  & \ddots    & \ddots & \ddots & \vdots  \\
    0       & \dots  & 0      & \lambda_{n-1} & z^{-1}t_{(n-1)n}\\
    0       & \dots  & 0      & 0     & \lambda_n
  \end{bmatrix},
  \]
  Elegimos ahora $z$ suficientemente grande de manera que
  \[
  \max_{1\le j \le n} \sum_{i=1}^{j-1} z^{i-j} |t_{ij}| < \varepsilon,
  \]
  y entonces
  \[
  \| D_z  U A \bar U^T D_z^{-1} \|_1
  = \max_{1\le j \le n} |\lambda_j| + \sum_{i=1}^{j-1} z^{i-j} |t_{ij}|
  < \rho(A) + \varepsilon.
  \]
  Faltaría ver que $\| D_z  U A \bar U^T D_z^{-1} \|_1$ es una norma inducida de $A$. Definimos entonces la norma $\| x \| = \| D_z U x \|_1$ para $x \in \mathbb{R}^n$ (que resulta una norma, verificar), y observemos que
  \[
  \| A \| = \max_{\substack{x\in\mathbb{R}^n\\x\neq 0}} \frac{\|Ax\|}{\|x\|}
  = \max_{\substack{x\in\mathbb{R}^n\\x\neq 0}} \frac{\|D_z U Ax\|_1}{\|\underbrace{D_z U x}_{y}\|_1}
  = \max_{\substack{y\in\mathbb{R}^n\\y\neq 0}} \frac{\|D_z U A (D_z U)^{-1} y\|_1}{\|y\|_1}
  = \| D_z U A \bar U^T D_z^{-1} \|_1 .
  \]
\end{proof}

La importancia del \emph{radio espectral} queda plasmada en el siguiente teorema.

\begin{thm}\label{T:equivalencia-radio-espectral}
  Sea $A\in\Rnn$, entonces las siguientes afirmaciones son equivalentes:
  \begin{enumerate}[label=\textup{(\alph*)}]
  \item $\rho(A) < 1$.
  \item $\|A \|<1$ para alguna norma inducida.
  \item $\displaystyle \lim_{k \to\infty} \|A^k\| = 0$, para toda norma inducida $\| \cdot \|$.
  \item $\displaystyle \lim_{k\to\infty} A^k x = 0$, para todo vector $x \in \mathbb{R}^n$.
  \end{enumerate}
\end{thm}

\begin{proof}
  $(1) \Rightarrow (2)$ es inmediato por el Teorema~\ref{T:radio-espectral}.

  $(2) \Rightarrow (3)$. Supongamos que $\|A\|_* < 1$ para alguna norma inducida $\| \cdot \|_*$. Sea $\| \cdot \|$ una norma inducida cualquiera. Entonces, como todas las normas son equivalentes, por ser el espacio de las matrices de dimensión finita, existe $c > 0$ tal que
  \[
  \| A^k \| \le c \| A^k \|_*, \qquad \forall k \in \mathbb{N}_0.
  \]
  Por el Teorema~\ref{T:consistencia} de consistencia, resulta que $\|A^k \|_* \le \|A\|_*^k$ y luego
  \[
  \| A^k \| \le c \| A^k \|_* \le c \|A\|_*^k \to 0 , \quad\text{cuando }k\to\infty.
  \]

  $(3) \Rightarrow (4)$. Sea $x \in \Rn$ y sea $\| \cdot \|$ una norma en $\Rn$ y también su norma inducida en $\Rnn$. Entonces, por el Teorema~\ref{T:consistencia} de consistencia
  \[
  \| A^k x - 0\| = \| A^k x \| \le \| A^k \| \| x \|  \to 0 , \quad\text{cuando }k\to\infty.
  \]

  $(4) \Rightarrow (1)$. Sea $\lambda \in \sigma(A)$ y sea $x$ un autovector asociado a $\lambda$, luego
  \[
  Ax = \lambda x, \qquad \text{y}\qquad
  A^k x = \lambda^k x.
  \]
  Así,
  \[
  \| A^k x \| = |\lambda|^k \| x\|
  \]
  para cualquier norma $\| \cdot \| $ de $\mathbb{R}^n$. Por (4), $\| A^k x \| \to 0$ y por ende $|\lambda|^k \| x\| \to 0$. Como $x$ es un autovector, es no nulo, y por lo tanto $|\lambda|^k \to 0$ cuando $k\to\infty$. Por lo tanto $|\lambda| < 1$. Como $\lambda\in\sigma(A)$ era arbitrario, resulta $\rho(A) = \max_{\lambda \in \sigma(A)} |\lambda| < 1$.
\end{proof}


\begin{thm} Sea $A\in\Rnn$. Entonces, cualquiera sea la norma inducida $\| \cdot \|$, resulta
  \[
  \rho(A) = \lim_{k \to \infty} \| A^k \|^{1/k}.
  \]
\end{thm}

\begin{proof}
  Veamos primero que $\rho(A)\le \| A^k \|^{1/k}$. Sea $\lambda \in \sigma(A)$ y sea $x$ un autovector asociado a $\lambda$. Entonces
  \begin{align*}
    Ax &= \lambda x \quad\Longrightarrow\quad
    A^k x = \lambda ^k x \quad\Longrightarrow\quad
    \| A^k x \| = |\lambda|^k \| x \| \\
    \Longrightarrow\quad |\lambda|^k &= \frac{\| A^k x \|}{\| x \|} \le \| A^k \|
    \quad\Longrightarrow\quad |\lambda| \le  \| A^k \|^{1/k}
    \quad\Longrightarrow\quad \rho(A)  \le  \| A^k \|^{1/k}.
  \end{align*}

  Veamos ahora que el límite de $\| A^k \|^{1/k}$ es $\rho(A)$.
  Sea $\varepsilon > 0$ y definimos
  \[
  \tilde A = \frac{1}{\rho(A) + \varepsilon} A \quad\Longrightarrow\quad
  \rho(\tilde A) = \frac{1}{\rho(A) + \varepsilon} \rho(A) = \frac{\rho(A)}{\rho(A)+\varepsilon} < 1.
  \]
  Por el Teorema anterior, $\displaystyle \lim_{k\to\infty} \|\tilde A^k \| = 0$ y
  por lo tanto, existe $k_0 \in \mathbb{N}$ tal que
  \[
  \| \tilde A^k \| = \frac{\|A^k\|}{(\rho(A)+\varepsilon)^k} < 1 , \quad\forall k \ge k_0.
  \]
  Es decir que dado $\varepsilon > 0$, existe $k_0\in\mathbb{N}$ tal que
  \[
  \rho(A) \le \| A^k \|^{1/k} \le \rho(A) + \varepsilon,  \quad\forall k \ge k_0.
  \]
  O sea, $\lim_{k\to\infty} \| A^k \|^{1/k} = \rho(A)$.
\end{proof}

Cuando $\rho(A)<1$, la serie geométrica de matrices converge y da la inversa de
$I-A$.

\begin{thm}[Serie de Neumann]\label{T:neumann}
Sea $A\in\Rnn$ con $\rho(A)<1$. Entonces $I-A$ es invertible y
\[
(I-A)^{-1}=\sum_{k=0}^{\infty}A^{k}=I+A+A^{2}+A^{3}+\cdots ,
\]
con convergencia en cualquier norma. Si además $\|A\|<1$ para alguna norma
inducida, entonces
\[
\bigl\|(I-A)^{-1}\bigr\|\ \le\ \frac{1}{1-\|A\|} ,
\qquad
\Bigl\|(I-A)^{-1}-\sum_{k=0}^{m}A^{k}\Bigr\|\ \le\ \frac{\|A\|^{m+1}}{1-\|A\|} .
\]
\end{thm}

\begin{proof}
Si $I-A$ no fuera invertible, existiría $x\neq0$ con $(I-A)x=0$, es decir
$Ax=x$: entonces $1$ sería autovalor de $A$ y $\rho(A)\ge1$, en contra de la
hipótesis. Sea $S_m=\sum_{k=0}^{m}A^{k}$. La identidad telescópica
\[
(I-A)\,S_m = I-A^{m+1}
\]
se verifica desarrollando el producto. Por el
Teorema~\ref{T:equivalencia-radio-espectral}, $\rho(A)<1$ implica
$\|A^{m+1}\|\to0$, así que $(I-A)S_m\to I$ y, multiplicando por
$(I-A)^{-1}$, $S_m\to(I-A)^{-1}$.

Para las cotas, si $\|\cdot\|$ es inducida y $\|A\|<1$, la submultiplicatividad
da $\|A^k\|\le\|A\|^k$ y entonces
$\|S_m\|\le\sum_{k=0}^m\|A\|^k\le 1/(1-\|A\|)$; pasando al límite se obtiene
la primera. La segunda sale de
$(I-A)^{-1}-S_m=\sum_{k>m}A^{k}=A^{m+1}(I-A)^{-1}$ y de la primera.
\end{proof}

\begin{remark}[De dónde viene esto]\label{obs:neumann-iterativos}
La serie de Neumann es la versión matricial de $\frac1{1-a}=\sum a^k$ para
$|a|<1$, y es el esqueleto de los \emph{métodos iterativos} para resolver
$Ax=b$. La idea es partir $A=M-N$ con $M$ fácil de invertir ---diagonal, o
triangular, como las que salieron de la eliminación en el
Capítulo~\ref{cap:LU}--- y escribir $Ax=b$ como
\[
x = M^{-1}Nx + M^{-1}b =: Tx + c .
\]
La iteración $x^{(k+1)}=Tx^{(k)}+c$ converge, para cualquier punto de partida,
exactamente cuando $\rho(T)<1$, y entonces el límite es
$(I-T)^{-1}c$, que es la solución. El radio espectral de la \emph{matriz de
iteración} $T$ es, entonces, lo único que hay que mirar para saber si un
método iterativo sirve, y cuán rápido converge: el error se multiplica,
aproximadamente, por $\rho(T)$ en cada paso.
\end{remark}

% ===================================================================
\section{Número de condición y error relativo}

Dada $A\in\Rnn$ y $b \in \mathbb{R}^n$, consideremos el problema de hallar $x \in \mathbb{R}^n$ tal que $Ax = b$. Supongamos que $\tilde x \in\mathbb{R}^n$ es una solución aproximada del sistema. Considerando una norma dada $\| \cdot \|$ de $\mathbb{R}^n$, definimos:
\begin{align*}
  \text{error} &=  x - \tilde x , \\
  \text{error relativo} &= \frac{ \| x - \tilde x \| }{ \| x \| },\\
  \text{residuo} &= A\tilde x - b = \tilde b - b, \quad\text{si $\tilde b = A \tilde x$},\\
  \text{residuo relativo} &= \frac{ \|\tilde b - b\|}{\|b\|}.
\end{align*}

Queremos hallar una estimación para el error relativo $\frac{ \| x - \tilde x \| }{ \| x \| }$ en términos del residuo relativo $ \frac{ \|\tilde b - b\|}{\|b\|}$.
Observemos lo siguiente
\[
\| \tilde x - x \| = \| A\inv \tilde b - A\inv b \| \le \| A\inv \| \| \tilde b - b \|,
\]
Además,
\[
\| A x \| = \| b \| \quad\Longrightarrow\quad \| b \| \le \| A \| \| x \|.
\]
Luego,
\begin{equation}\label{eq:err_rel1}
\frac{\| \tilde x - x \| }{\| x \|}
\le \| A \| \| A\inv\| \frac{\| \tilde b - b \|}{\| b \|}.
\end{equation}

Análogamente
\[
\| \tilde b - b \| = \| A \tilde x - A x \| \le \| A \| \| \tilde x - x \|,
\]
y
\[
\| x \| = \| A\inv b \| \le \| A\inv \| \| b \|.
\]
Por lo tanto
\begin{equation}\label{eq:err_rel2}
  \frac{1}{\| A \| \| A\inv\|} \frac{\| \tilde b - b \|}{\| b \|}
  \le \frac{\| \tilde x - x \| }{\| x \|}.
\end{equation}

De~\eqref{eq:err_rel1}--\eqref{eq:err_rel2} obtenemos que
\begin{equation}\label{eq:err_rel}
  \frac{1}{\| A \| \| A\inv\|} \frac{\| \tilde b - b \|}{\| b \|}
  \le \frac{\| \tilde x - x \| }{\| x \|}
  \le \| A \| \| A\inv\| \frac{\| \tilde b - b \|}{\| b \|}.
\end{equation}

\begin{defn}
  Dada una norma inducida $\| \cdot \|$ en el espacio de las matrices de $n\times n$, se define el \textbf{número de condición} para esa norma como
  \[
  \kappa (A ) =  \| A \| \, \| A\inv\|, \qquad A \in \Rnn \text{ invertible}.
  \]
\end{defn}

\begin{lema}\label{L:error-vs-residuo}
Si $x \in \Rn \setminus\{0\}$ es la solución de $Ax = b$, entonces, cualquiera sea $\tilde x \in \Rn$, se cumple que
\[
  \frac{1}{\kappa(A)} \frac{\| \tilde b - b \|}{\| b \|}
  \le \frac{\| \tilde x - x \| }{\| x \|}
  \le \kappa(A) \frac{\| \tilde b - b \|}{\| b \|},
  \]
  para $\tilde b = A \tilde x$.
\end{lema}



\begin{remark} El número de condición siempre es mayor o igual a uno:
  \[
  \kappa(A) =  \| A \| \| A\inv\| \ge \| A A\inv \| = \| I \| = 1.
  \]
  Cuando $\kappa(A)$ es pequeño, o cercano a uno, $\kappa(A) \cong 1$ resulta
  \[
  \frac{\| \tilde x - x \| }{\| x \|}
  \cong \frac{\| \tilde b - b \|}{\| b \|}
  \]
  Es decir, pequeñas perturbaciones en el lado derecho de la ecuación, producen pequeños cambios en las soluciones. En este caso decimos que el problema $Ax=b$ está \emph{bien condicionado}.
\end{remark}

\begin{remark}
  Dada la solución $x$ de $Ax = b$, si $\tilde x$ es una aproximación de $x$, entonces
  \[
  \underbrace{\frac{\| \tilde x - x \| }{\| x \|}}_{\text{desconocido}} \le \kappa(A)
  \underbrace{\frac{\| A \tilde x - b \|}{\| b \|}}_{\text{calculable}}.
  \]
  El lado derecho de la desigualdad anterior se utiliza en la implementación de criterios de parada en algunos métodos iterativos.
\end{remark}

\begin{remark} Observemos que multiplicar una matriz por una constante arbitraria no nula, no afecta el número de condición:
  \[
  \kappa(\alpha A) = \| \alpha A \| \| (\alpha A)\inv \|
  = |\alpha| \| A \| \frac{1}{|\alpha|} \| A\inv\| = \kappa(A).
  \]
\end{remark}

Hay una relación estrecha entre el número de condición y el espectro cuando se
usa la norma $\|\cdot\|_2$ y la matriz es simétrica.

\begin{lema}\label{L:condicion-simetrica}
  Si $A\in\Rnn$ es simétrica e invertible, entonces
  \[
  \kappa_2(A)=\frac{|\lambda|_{\max}}{|\lambda|_{\min}} ,
  \]
  donde $|\lambda|_{\max}$ y $|\lambda|_{\min}$ son el mayor y el menor
  \emph{módulo} de los autovalores de $A$.
\end{lema}

\begin{proof}
  Como $A$ es simétrica, $\|A\|_2=\rho(A)=\max\{|\lambda|:\lambda\in\sigma(A)\}
  =|\lambda|_{\max}$. Como además $A$ es invertible, $A\inv$ es simétrica y sus
  autovalores son los $\lambda\inv$ con $\lambda\in\sigma(A)$
  (Corolario~\ref{cor:traza-det}), así que
  \begin{align*}
  \|A\inv\|_2 &= \rho(A\inv)
  = \max\bigl\{|\mu| : \mu\in\sigma(A\inv)\bigr\}
  = \max\bigl\{|\lambda|\inv : \lambda\in\sigma(A)\bigr\}\\
  &= \frac{1}{\min\bigl\{|\lambda| : \lambda\in\sigma(A)\bigr\}}
  = \frac{1}{|\lambda|_{\min}} ,
  \end{align*}
  donde la anteúltima igualdad dice simplemente que el máximo de los inversos
  de números positivos es el inverso del mínimo. Multiplicando las dos normas
  queda el enunciado.
\end{proof}

Para una matriz cualquiera ---sin ninguna hipótesis--- la respuesta la da la
SVD, y es todavía más simple.

\begin{lema}[El número de condición son los valores singulares]\label{L:condicion-svd}
Sea $A\in\Rnn$ invertible, con valores singulares
$\sigma_1\ge\cdots\ge\sigma_n>0$. Entonces
\[
\boxed{\ \kappa_2(A)=\frac{\sigma_1}{\sigma_n}\ }
\]
\end{lema}

\begin{proof}
Por la Definición~\ref{def:norma-espectral}, $\|A\|_2=\sigma_1$. Si
$A=U\Sigma V^T$ es una SVD de $A$ con $\Sigma=\diag(\sigma_1,\dots,\sigma_n)$
invertible, entonces
\[
A\inv=V\Sigma^{-1}U^T ,
\]
y como $U$ y $V$ son ortogonales eso es una SVD de $A\inv$, salvo por el orden
de la diagonal: los valores singulares de $A\inv$ son
$\sigma_n^{-1}\ge\cdots\ge\sigma_1^{-1}$. El mayor es $\sigma_n^{-1}$, de
donde $\|A\inv\|_2=1/\sigma_n$.
\end{proof}

\begin{remark}[Qué mide, en una frase]\label{obs:condicion-elipsoide}
El Teorema~\ref{teo:elipsoide} dice que $A$ transforma la esfera unitaria en
un elipsoide de semiejes $\sigma_1,\dots,\sigma_n$. El número de condición es
el cociente entre el semieje más largo y el más corto: cuán \emph{achatado}
queda ese elipsoide. Una matriz bien condicionada manda la esfera a algo
todavía parecido a una esfera; una mal condicionada la aplasta casi hasta un
subespacio, y por eso un cambio chico del lado derecho produce un cambio
enorme en la solución. De paso se ve que $\kappa_2(A)=1$ exactamente cuando
todos los valores singulares coinciden, es decir cuando $A$ es un múltiplo de
una matriz ortogonal: \emph{las matrices ortogonales son las mejor
condicionadas que hay}, y es la razón por la que $QR$ y las reflexiones de
Householder del Capítulo~\ref{cap:ortogonalidad} son numéricamente preferibles
a la eliminación gaussiana.
\end{remark}





% ===================================================================
\section*{Ejercicios}

\begin{enumerate}

% ---------- normas de matrices: definiciones y propiedades ----------
\item\label{ejer-normaiii} Mostrar que si $\|\cdot\|$ es una norma en $\RR^{np}$ entonces $ \vertiii{\cdot}$ en $\RR^{n\times p}$ definido como
\[\vertiii{A} = \|\textup{vec}(A)\|\]
es una norma.

\item\label{ej-matrizortogonal} Sea $A\in\RR^{m\times n}$.
\begin{enumerate}[label=\textup{(\alph*)}]
    \item Mostrar que $\|A\|_F^2 = \sum_{j=1}^n \|\col_j(A)\|_2^2$ y también
    $\|A\|_F^2 = \sum_{i=1}^m \|\fila_i(A)\|_2^2$: la norma de Frobenius al
    cuadrado es tanto la suma de los cuadrados de las longitudes de las
    columnas como la de las filas.
    \item Mostrar que si $Q$ es una matriz ortogonal entonces $\|A\|_F = \|QA\|_F = \|AQ\|_F$.
\end{enumerate}

\item Decimos que una norma de matrices está \textbf{normalizada} si la matriz
identidad de cualquier orden tiene norma $1$. (Cuidado con el nombre: no tiene
nada que ver con las matrices unitarias del Capítulo~\ref{cap:schur} ni con
las normas ortogonalmente invariantes de la Sección~\ref{sec:tres-normas}; de
hecho $\|\cdot\|_F$ es ortogonalmente invariante y no está normalizada.)
\begin{enumerate}[label=\textup{(\alph*)}]
    \item Mostrar que toda norma inducida está normalizada.
    \item ¿Es cierto que toda norma matricial debe estar normalizada?
\item Considere en $\RR^{n\times p}$
\[\|A\| = \|\textup{vec}(A)\|_1.\]
¿Está normalizada? ¿Es una norma matricial?
\end{enumerate}

\item\label{ej-equivalencianormainducida} Probar que 
\[ \max_{\bec x \neq \bec 0}\frac{\|A\bec x\|}{\|\bec x\|} = \max_{\|\bec x \| = 1}{\|A\bec x\|} = \max_{\|\bec x \| \leq  1}{\|A\bec x\|}   \]

\item Para una matriz $A$ de columnas $\{\bec a_1, \dots, \bec a_n\}$ definimos
\[\|A\|_{2,1} = \sum_{j=1}^n \|\bec a_j\|_2.\]
Mostrar que es una norma. ¿Es una norma matricial?

\item\label{ej-medica} Mostrar que la siguiente modificación de la norma del
máximo la vuelve una norma matricial en $\RR^{m\times n}$:
\[\|A\| = \sqrt{mn}\,\|\textup{vec}(A)\|_{\infty} .\]

\item\label{ej-traspuesta} Probar que $\|A\|_F=\|A^T\|_F$ y
$\|A\|_2=\|A^T\|_2$, pero que en cambio $\|A\|_1=\|A^T\|_\infty$. Dar una
matriz con $\|A\|_1\neq\|A^T\|_1$.

\item\label{ej-mixtas} \emph{(Cotas mixtas.)} Usando que
$\fila_i(AB)=\fila_i(A)B$ ---la fórmula~\eqref{eq:filas} del
Capítulo~\ref{cap:algebra-lineal}--- probar que
$\|AB\|_F\le\|A\|_F\,\|B\|_2$, que es la versión por filas de la
Observación~\ref{obs:frob-por-columnas}. ¿Por qué estas dos cotas son mejores
que la submultiplicatividad $\|AB\|_F\le\|A\|_F\|B\|_F$ de la
Proposición~\ref{prop:frobenius-matricial}?

% ---------- cuentas a mano ----------
\item\label{ej-calcular-normas} \emph{(A mano.)} Para cada una de las matrices
\[
A_1=\begin{bmatrix}3&0\\ 4&5\end{bmatrix},\qquad
A_2=\begin{bmatrix}0&2\\ 0&0\end{bmatrix},\qquad
A_3=\begin{bmatrix}1&1\\ 1&1\end{bmatrix},
\]
\[
A_4=\begin{bmatrix}2&-1\\ -1&2\end{bmatrix},\qquad
A_5=\begin{bmatrix}0&-1\\ 1&0\end{bmatrix},\qquad
A_6=\diag(3,-2,1),
\]
calcular, haciendo las cuentas:
$\|A\|_1$, $\|A\|_\infty$, $\|A\|_F$, $\|A\|_2$, $\|A\|_*$ y $\rho(A)$.
\emph{Sugerencias:} $\|A\|_1$ y $\|A\|_\infty$ salen del
Teorema~\ref{teo-l1l2linfty} mirando sumas de columnas y de filas; para
$\|A\|_2$ y $\|A\|_*$ hacen falta los valores singulares, que acá se consiguen
sin esfuerzo en cada caso: $A_1$ tiene su SVD calculada en el
Ejemplo~\ref{ej:svd-completo}; $A_2$ y $A_3$ tienen rango uno;
$A_4$, $A_6$ son simétricas (Ejercicio~\ref{ej-simetrica-rho}); y $A_5$ es
ortogonal. Armar una tabla con los seis resultados en filas.

\item\label{ej-orden-normas} Con los números del ejercicio anterior:
\begin{enumerate}[label=\textup{(\alph*)}]
\item Verificar en cada caso las desigualdades
$\rho(A)\le\|A\|_2\le\|A\|_F\le\sqrt r\,\|A\|_2$, donde $r=\rango(A)$.
\item Probar las dos últimas en general, y caracterizar las matrices en las
que vale cada igualdad. \emph{Sugerencia:} las tres normas se leen de los
$\sigma_k$ (Sección~\ref{sec:tres-normas}); la primera igualdad pide
$\sigma_1^2=\sigma_1^2+\cdots+\sigma_r^2$ y la segunda, que todos los
$\sigma_k$ sean iguales. ¿Cuáles de las seis matrices de arriba realizan cada
igualdad?
\item Deducir de la misma manera que $\|A\|_F\le\|A\|_*\le\sqrt r\,\|A\|_F$.
\item ¿En cuáles de las seis es $\rho(A)<\|A\|_2$ estrictamente? ¿Qué tienen
en común aquellas en las que vale la igualdad?
\end{enumerate}

\item\label{ej-rango-uno-normas} \emph{(Rango uno.)} Sea $A=\bec u\bec v^T$
con $\bec u\in\RR^m$ y $\bec v\in\RR^n$ no nulos.
\begin{enumerate}[label=\textup{(\alph*)}]
\item Probar que $A$ tiene un solo valor singular no nulo y que vale
$\|\bec u\|_2\|\bec v\|_2$. \emph{Sugerencia:} $A^TA=\|\bec u\|_2^2\,\bec v\bec v^T$.
\item Deducir que $\|A\|_2=\|A\|_F=\|A\|_*=\|\bec u\|_2\|\bec v\|_2$.
\item Probar que $\|A\|_1=\|\bec u\|_1\max_j|v_j|$ y
$\|A\|_\infty=\|\bec v\|_1\max_i|u_i|$.
\item Verificar los cuatro valores a mano con $\bec u=(1,2,-2)^T$ y
$\bec v=(3,0,4)^T$. (Deberían dar $15$, $15$, $15$, $20$ y $14$.)
\end{enumerate}

\item\label{ej-identidad-unos} Calcular $\|\cdot\|_1$, $\|\cdot\|_\infty$,
$\|\cdot\|_F$, $\|\cdot\|_2$ y $\|\cdot\|_*$ de la identidad $I_n$ y de la
matriz $\mathbf 1\mathbf 1^T\in\RR^{n\times n}$ de todos unos. Las dos
matrices tienen la misma norma $\|\cdot\|_2$; ¿cuál de las normas distingue
mejor el rango, y por qué?

% ---------- del Teorema 12.5 y del radio espectral ----------
\item\label{ej-linfty} \emph{(La parte que falta del
Teorema~\ref{teo-l1l2linfty}.)} Probar la parte~\ref{item-linfty}, es decir
que $\|A\|_\infty=\max_i\|\fila_i(A)\|_1$, siguiendo el mismo esquema que la
demostración de la parte~\ref{item-l1} pero mirando por filas.
\begin{enumerate}[label=\textup{(\alph*)}]
\item Probar que si $\|\bec x\|_\infty\le1$ entonces
$|(A\bec x)_i|=|\fila_i(A)\,\bec x|\le\|\fila_i(A)\|_1$, y deducir
$\|A\|_\infty\le\max_i\|\fila_i(A)\|_1$.
\item Para la otra desigualdad, sea $i_0$ la fila de mayor norma $\ell_1$ y
tomar $\bec x$ con $x_j=\sgn(a_{i_0j})$ (y $x_j=1$ si
$a_{i_0j}=0$). Verificar que $\|\bec x\|_\infty=1$ y que
$(A\bec x)_{i_0}=\|\fila_{i_0}(A)\|_1$.
\item ¿Por qué acá el vector que realiza el máximo tiene todas sus
coordenadas no nulas, mientras que en la parte~\ref{item-l1} alcanzaba con un
vector de la base canónica?
\end{enumerate}

\item\label{ej-simetrica-rho} Probar que si $A\in\Rnn$ es simétrica entonces
$\|A\|_2=\rho(A)$ ---la igualdad que se usa sin demostrar al final de la
Sección~\ref{sec:normas-inducidas}--- y que sus valores singulares son los
módulos de sus autovalores. Dar una matriz no simétrica con
$\rho(A)<\|A\|_2$ y otra, también no simétrica, con $\rho(A)=\|A\|_2$.

\item\label{ej-rho-cota} Probar que $\rho(A)\le\|A\|$ para toda norma inducida
$\|\cdot\|$ y toda $A$ cuadrada. \emph{Sugerencia:} si $A\bec x=\lambda\bec x$
con $\bec x\neq\bec 0$, dividir $\|A\bec x\|$ por $\|\bec x\|$. Dar después
una matriz con $\rho(A)=0$ y $\|A\|_2$ tan grande como se quiera, y concluir
que $\rho$ \emph{no} es una norma.

\item\label{ej-equivalencias} Del Teorema~\ref{T:equivalencia-radio-espectral}
demostrar las implicaciones $(c)\Rightarrow(d)$ y $(d)\Rightarrow(a)$, que no
usan nada más que la submultiplicatividad y la definición de autovector.

\item\label{ej-potencias} Verificar $\rho(A)=\lim_k\|A^k\|^{1/k}$ en
$A=\begin{bsmallmatrix}0&t\\ 0&0\end{bsmallmatrix}$ con $t>0$ y en una matriz
diagonalizable a elección, calculando $\|A^k\|^{1/k}$ para
$k=1,2,5,10,50$. ¿Por qué el límite se alcanza enseguida en un caso y es lento
en el otro?

% ---------- más finos ----------
\item\label{ej-2-1-infty} Probar que $\|A\|_2^2\le\|A\|_1\,\|A\|_\infty$.
\emph{Sugerencia:} $\|A\|_2^2=\|A^TA\|_2\le\|A^TA\|_\infty$ y
$\|A^T\|_\infty=\|A\|_1$ (Ejercicio~\ref{ej-traspuesta}). Verificarlo en
$A_1=\begin{bsmallmatrix}3&0\\ 4&5\end{bsmallmatrix}$ del
Ejercicio~\ref{ej-calcular-normas}.

\item\label{ej-nuclear} Probar que $\|A\|_*=\sum_i\sigma_i(A)$ es una norma.
Lo único que cuesta es la desigualdad triangular; una manera es usar la
caracterización
\[
\|A\|_*=\max\bigl\{\tr(Q^TA)\ :\ Q\in\R^{m\times n},\ \|Q\|_2\le1\bigr\} ,
\]
que conviene probar primero con la SVD de $A$: el máximo se alcanza en
$Q=U_rV_r^T$.

\item\label{ej-condicion} \emph{(Número de condición.)}
\begin{enumerate}[label=\textup{(\alph*)}]
\item Probar que $\kappa(A)\ge1$ para toda norma inducida y que
$\kappa(\alpha A)=\kappa(A)$ para todo $\alpha\neq0$.
\item Calcular $\kappa_2(A)$ para
$A=\begin{bsmallmatrix}1&1\\ 1&1+\varepsilon\end{bsmallmatrix}$ y estudiar qué
pasa cuando $\varepsilon\to0$. (Para $\varepsilon$ chico da del orden de
$4/\varepsilon$.)
\item Resolver $Ax=b$ con $b=(2,2)^T$ y con $b=(2,2+\delta)^T$ y comparar el
cambio relativo de la solución con el del lado derecho, contrastándolo con la
cota del Lema~\ref{L:error-vs-residuo}.
\end{enumerate}

\item Probar que si $A = \bec u \bec v^T$ y $B= \bec v \bec w^T$ con
$\bec u,\bec v,\bec w$ vectores no nulos, entonces se satisface que
$\|AB\|_F^2=\|A\|_F^2\|B\|_F^2$.

La implicación contraria es una conjetura: las únicas matrices sin elementos
nulos, de rango uno, que satisfacen $\|AB\|_F=\|A\|_F\,\|B\|_F$ son de la
forma $A=\bec u\bec v^T$ y $B=\bec v\bec w^T$ con un vector compartido
$\bec v$.
\end{enumerate}
